% ======================================================================
\section{From the light-cone Hamiltonian to its eikonal form}\label{app:hamiltonian}
% ======================================================================

This appendix gives the steps between the light-cone Hamiltonian~\eqref{eq:HLC} and the eikonal Hamiltonian of the soft modes used in the main text. It follows Appendices~A and~B of~\cite{Lublinsky:2016meo}, restricted to pure Yang--Mills theory.

\subsection{Expansion in the coupling and quantisation}

In Eq.~\eqref{eq:PiF} all fields are taken at the same light-cone time. Inserting $D^{ab}_i=\delta^{ab}\partial_i-gf^{abc}A^c_i$ gives $\Pi^a=-\partial_iA^a_i+gf^{abc}\frac{1}{\partial^+}\big(A^c_i\,\partial^+A^b_i\big)$, and $\frac12\Pi^a\Pi^a$ contains the terms
\begin{align}
\frac12(\partial_iA^a_i)^2\,,\qquad -gf^{abc}(\partial_iA^a_i)\frac{1}{\partial^+}(A^b_j\partial^+A^c_j)\,,\qquad \frac{g^2}{2}\Big[f^{abc}\frac{1}{\partial^+}(A^b_i\partial^+A^c_i)\Big]^2\,.
\end{align} The magnetic part $\frac14F^a_{ij}F^a_{ij}$ contains $\frac12(\partial_iA^a_j)^2-\frac12(\partial_iA^a_j)(\partial_jA^a_i)$, the cubic term $-gf^{abc}(\partial_iA^a_j)A^b_iA^c_j$ and the quartic term $\frac{g^2}{4}f^{abc}f^{ade}A^b_iA^c_jA^d_iA^e_j$. After an integration by parts in the transverse coordinates, $(\partial_iA^a_j)(\partial_jA^a_i)\to(\partial_iA^a_i)^2$, the terms quadratic in the field combine into the free part. Expanding in powers of $g$ gives $H=H_0+H_{\rm int}$. The free part is
\begin{align}\label{eq:H0x}
H_0=\int dx^-\,d^2x\ \frac12\,\big(\partial_iA^a_j\big)\big(\partial_iA^a_j\big)\,,
\end{align}
and the interaction is
\begin{align}\label{eq:Hint}
H_{\rm int}=\int dx^-\,d^2x\,\Big[&-gf^{abc}\,\big(\partial_iA^a_j\big)A^b_iA^c_j+\frac{g^2}{4}\,f^{abc}f^{ade}\,A^b_iA^c_jA^d_iA^e_j\notag\\
&-gf^{abc}\,\big(\partial_iA^a_i\big)\frac{1}{\partial^+}\big(A^b_j\,\partial^+A^c_j\big)\notag\\
&+\frac{g^2}{2}\,f^{abc}f^{ade}\,\frac{1}{\partial^+}\big(A^b_i\,\partial^+A^c_i\big)\frac{1}{\partial^+}\big(A^d_j\,\partial^+A^e_j\big)\Big]\,.
\end{align}
The first two terms are the three- and four-gluon vertices that originate from the magnetic part. The third and fourth terms originate from the electric part; they contain the inverse of $\partial^+$ and are the instantaneous interactions characteristic of light-cone quantisation. Eq.~\eqref{eq:Hint} follows from Eqs.~\eqref{eq:HLC} and~\eqref{eq:PiF} by direct expansion, with integrations by parts in the transverse coordinates in the free part.

The theory is quantised by expanding the transverse field at $x^+=0$ in creation and annihilation operators,
\begin{align}\label{eq:modeexp}
A^a_i(x)=\int_0^\infty\frac{dk^+}{2\pi}\int\frac{d^2k}{(2\pi)^2}\,\frac{1}{\sqrt{2k^+}}\Big(a^a_i(k^+,\v k)\,e^{-ik\cdot x}+a^{\dagger a}_i(k^+,\v k)\,e^{ik\cdot x}\Big)\,,
\end{align}
with $k\cdot x=k^+x^--\v k\cdot\v x$ at $x^+=0$ and the commutator~\eqref{eq:aadagger}. Inserted into Eq.~\eqref{eq:H0x} this gives, after normal ordering,
\begin{align}\label{eq:H0}
H_0=\int_0^\infty\frac{dk^+}{2\pi}\int\frac{d^2k}{(2\pi)^2}\,\frac{\v k^2}{2k^+}\,a^{\dagger a}_i(k^+,\v k)\,a^a_i(k^+,\v k)\,.
\end{align}
The eigenstates of $H_0$ are the free Fock states, obtained by acting with creation operators on the perturbative vacuum; the energy of a state is the sum of the energies~\eqref{eq:energy} of its gluons. The interaction $H_{\rm int}$ generates transitions between Fock states. Its form in terms of creation and annihilation operators is needed only after the separation of the modes, which we discuss next.


\subsection{Soft and valence modes; the eikonal terms}

With the decomposition~\eqref{eq:Asplit} and the eikonal approximation described in Section~\ref{sec:BO}, the interaction~\eqref{eq:Hint} reduces to the four terms of Eq.~\eqref{eq:Hdecomp},
\begin{align}\label{eq:Hpieces}
H_g&=-g\int dx^-d^2x\,\big(\partial_i\underline A^a_i\big)\,f^{abc}\,\frac{1}{\partial^+}\big(\overline A^b_j\,\partial^+\overline A^c_j\big)\,,\notag\\
H_{ggg}&=-gf^{abc}\int dx^-d^2x\,\Big[\big(\partial_i\underline A^a_j\big)\underline A^b_i\underline A^c_j+\big(\partial_i\underline A^a_i\big)\frac{1}{\partial^+}\big(\underline A^b_j\,\partial^+\underline A^c_j\big)\Big]\,,\notag\\
H_{gg\text{-inst}}&=g^2f^{abc}f^{ade}\int dx^-d^2x\,\frac{1}{\partial^+}\big(\underline A^b_i\,\partial^+\underline A^c_i\big)\,\frac{1}{\partial^+}\big(\overline A^d_j\,\partial^+\overline A^e_j\big)\,,\notag\\
H_{gggg}&=\frac{g^2}{4}f^{abc}f^{ade}\int dx^-d^2x\,\Big[\underline A^b_i\underline A^c_j\underline A^d_i\underline A^e_j+2\,\frac{1}{\partial^+}\big(\underline A^b_i\,\partial^+\underline A^c_i\big)\frac{1}{\partial^+}\big(\underline A^d_j\,\partial^+\underline A^e_j\big)\Big]\,.
\end{align}
$H_g$ is the coupling of one soft gluon to the valence charge; $H_{ggg}$ is the three-gluon vertex of the soft modes, including its instantaneous part; $H_{gg\text{-inst}}$ is the instantaneous interaction of two soft gluons with the valence charge; $H_{gggg}$ contains the four-gluon vertex of the soft modes and the instantaneous interaction among four soft gluons. $H_{gggg}$ is listed for completeness: it does not enter any of the coefficients needed in this paper (Section~\ref{sec:threeout}).

In terms of the soft creation and annihilation operators the free Hamiltonian of the soft modes and the instantaneous term that contains the valence charge are (the term $H_g$ is given in Eq.~\eqref{eq:Hg})
\begin{align}
H_0&=\int_\Lambda^\vee\frac{dk^+}{2\pi}\int\frac{d^2k}{(2\pi)^2}\,\frac{\v k^2}{2k^+}\,a^{\dagger a}_i(k^+,\v k)\,a^a_i(k^+,\v k)\,,\label{eq:H0soft}\\
H_{gg\text{-inst}}&=\int_\Lambda^\vee\frac{dp^+}{2\pi}\frac{dq^+}{2\pi}\int\frac{d^2p}{(2\pi)^2}\frac{d^2q}{(2\pi)^2}\,\frac{ig^2f^{abc}}{\sqrt{p^+q^+}}
\Bigg(\frac{q^+\,\rho^a(-\v p-\v q)}{2(p^++q^+)^2}\,a^{\dagger b}_i(p)\,a^{\dagger c}_i(q)-{\rm h.c.}\notag\\
&\hspace{6.2cm}-\frac{(p^++q^+)\,\rho^a(-\v p+\v q)}{2(p^+-q^+)^2}\,a^{\dagger b}_i(p)\,a^c_i(q)\Bigg)\,.\label{eq:Hinst}
\end{align}
The first two terms of $H_{gg\text{-inst}}$ create or annihilate a pair of soft gluons. The last term transfers the longitudinal momentum $p^+-q^+$ and the transverse momentum $\v p-\v q$ between a soft gluon and the valence modes. For $H_{ggg}$ we shall need only the matrix element between a one-gluon and a two-gluon state, which is given in Eq.~\eqref{eq:m3} below.

The forms~\eqref{eq:Hpieces} and~\eqref{eq:Hinst} are those of~\cite{Lublinsky:2016meo}. The four-gluon interaction $H_{gggg}$ is not used in this paper (Section~\ref{sec:threeout}); the sign of its instantaneous term is written as it follows from $\frac12\Pi^a\Pi^a$.
